% !TeX root = ../复分析笔记.tex
% 由主文档通过 \input 载入；请编译主文档。
\setcounter{chapter}{10} % 保留原书章号。
\chapter{共形映射与典型区域的变换}
\label{chap:ca-conformal}

\section{局部保角与整体双全纯}

把复杂区域变成圆盘或半平面，能够简化积分和边值问题。但一个可逆坐标变换若任意扭曲角度，可能破坏原来的解析结构；全纯函数的局部伸缩与旋转恰好给出合适的变换。

\begin{definition}[共形与双全纯]
若 $f$ 在 $a$ 附近全纯且 $f'(a)\ne0$，则称它在 $a$ \emph{共形（conformal）}，这里采用保持定向的保角约定。若 $f:D\to G$ 是一一对应的全纯映射，且反函数也全纯，则称它为\emph{双全纯映射（biholomorphic map）}，也称两个区域之间的共形映射。
\end{definition}

由第~\ref{chap:ca-counting} 章的结果，单射全纯映射的导数不为零，且其反函数自动全纯。因此实际构造时，证明全纯、单射和满射即可。只证明导数不为零，得到的是局部可逆和局部保角，不足以说明整个区域上的单射；$e^z$ 在整个平面上的周期性就是反例。

对经过 $a$ 的两条光滑曲线，像曲线的切向量均乘以 $f'(a)$，所以其有向夹角不变。长度通常会乘以 $|f'(a)|$，面积在一阶近似中会乘以 $|f'(a)|^2$。因此共形不意味着保持距离或面积，也不意味着把有限大小的图形变成相似图形。

\section{莫比乌斯变换}

\begin{definition}[莫比乌斯变换]
函数
\begin{equation}
 T(z)=\frac{az+b}{cz+d},\qquad a,b,c,d\in\mathbb C,\quad ad-bc\ne0,
 \label{eq:ca-mobius}
\end{equation}
称为\emph{莫比乌斯变换（Möbius transformation）}，又称分式线性变换。同乘所有系数以一个非零常数不改变变换。
\end{definition}

若 $c\ne0$，在扩充复平面上规定 $T(-d/c)=\infty$、$T(\infty)=a/c$；若 $c=0$，规定 $T(\infty)=\infty$。解出反函数
\[
 T^{-1}(w)=\frac{dw-b}{a-cw},
\]
可知 $T$ 在球面上是一一对应。普通点处的导数为
\[
 T'(z)=\frac{ad-bc}{(cz+d)^2}\ne0.
\]
在极点和无穷远点附近，用 $1/z$ 或 $1/w$ 作局部坐标，同样得到非零导数。条件 $ad-bc\ne0$ 正是排除分子、分母比例相同而退化为常函数的情形。

\begin{proposition}[广义圆的保持性]
莫比乌斯变换把广义圆映为广义圆。这里广义圆指普通圆，或补上无穷远点的直线。
\end{proposition}
\begin{proof}
圆或直线可统一写成
\[
 A|z|^2+Bz+\overline B\,\overline z+C=0,
 \quad A,C\in\mathbb R,
\]
并要求方程表示非退化的圆或直线。将莫比乌斯反函数代入，乘以分母的模平方后，仍得到同一类型的方程。由于变换在球面上双射，原来的圆不会退化为一个点或空集，所以像仍为广义圆。
\end{proof}

这个性质只确定边界可能变成圆或直线，不能自动确定区域映到哪一侧。常用办法是先确定边界，再测试一个内部点；若还需证明整个区域一一对应，应结合反函数或明确的不等式。

\begin{proposition}[三点确定变换]
给定扩充复平面上的两个有序三元组，且各组内三点互异，存在唯一莫比乌斯变换把第一组依次映到第二组。
\end{proposition}
\begin{proof}
对互异有限点 $z_1,z_2,z_3$，
\[
 S(z)=\frac{(z-z_1)(z_2-z_3)}{(z-z_3)(z_2-z_1)}
\]
依次把它们映到 $0,1,\infty$；有无穷远点时取相应极限即可。对目标三点构造 $S_2$，所求为 $S_2^{-1}\circ S_1$。若两个变换都满足条件，它们的复合之差归结为一个固定 $0,1,\infty$ 的莫比乌斯变换；固定无穷远使 $c=0$，固定零使 $b=0$，固定 $1$ 使比例系数为 $1$，故只能是恒等变换。
\end{proof}

\section{上半平面、右半平面与单位圆盘}

三点归一化可以用\emph{交比（cross ratio）}记录。对四个互异有限点，约定
\[
 [z,z_1;z_2,z_3]
 =\frac{(z-z_1)(z_2-z_3)}{(z-z_3)(z_2-z_1)};
\]
有一个点为无穷远时取相应极限。不同教材可能采用不同排列约定，使用时应先比较定义。若 $T(z)=(az+b)/(cz+d)$，则在各分母非零处
\[
 T(u)-T(v)=\frac{(ad-bc)(u-v)}{(cu+d)(cv+d)}.
\]
将此式代入交比，所有附加因子相消，得到
$[T(z),T(z_1);T(z_2),T(z_3)]=[z,z_1;z_2,z_3]$；涉及无穷远点时由连续性取极限。交比不变性因此把三点对应转化为求未知像点的一个分式方程。

\begin{example}[凯莱变换]
\label{ex:ca-cayley}
证明 $T(z)=(z-i)/(z+i)$ 把上半平面双全纯地映到单位圆盘，并将 $i$ 映到 $0$。
\end{example}
\begin{solve}
令 $z=x+iy$，则
\[
 |z+i|^2-|z-i|^2=4y.
\]
所以 $y>0$ 当且仅当 $|T(z)|<1$。反函数为 $z=i(1+w)/(1-w)$，当 $|w|<1$ 时
\[
 \operatorname{Im}z=\frac{1-|w|^2}{|1-w|^2}>0.
\]
两式给出满射与反函数的定义域，莫比乌斯性质给出单射和全纯性。实轴与无穷远点组成的边界映到单位圆，$i$ 映到零。
\end{solve}

类似地，$S(z)=(z-1)/(z+1)$ 将右半平面 $\operatorname{Re}z>0$ 映到单位圆盘，因为 $|z+1|^2-|z-1|^2=4\operatorname{Re}z$。半平面的方向不同，公式中的基准点也不同，不能不经检查就互换 $1$ 与 $i$。

\begin{proposition}[圆盘自同构]
\label{prop:ca-disk-auto}
对 $a\in\mathbb D$，函数
\[
 \varphi_a(z)=\frac{z-a}{1-\overline a z}
\]
是单位圆盘的双全纯自映射，将 $a$ 映到 $0$，反函数为 $\varphi_a^{-1}(w)=(w+a)/(1+\overline a w)$。单位圆盘的全部双全纯自映射恰为
\[
 e^{i\theta}\frac{z-a}{1-\overline a z},\qquad a\in\mathbb D,\quad\theta\in\mathbb R.
\]
\end{proposition}
\begin{proof}
直接展开有
\begin{equation}
 1-|\varphi_a(z)|^2
 =\frac{(1-|a|^2)(1-|z|^2)}{|1-\overline a z|^2}.
 \label{eq:ca-disk-auto}
\end{equation}
右侧对 $z\in\mathbb D$ 为正，反函数也有相同性能，故得到双射。对任意双全纯自映射 $F$，取 $a=F^{-1}(0)$，令 $H=F\circ\varphi_a^{-1}$，则 $H(0)=0$。施瓦茨引理分别应用于 $H$ 和 $H^{-1}$，得 $|H(z)|\le|z|\le|H(z)|$，所以处处取等号，$H(z)=e^{i\theta}z$，从而得到分类。
\end{proof}

\begin{theorem}[施瓦茨--皮克不等式，Schwarz--Pick inequality]
\label{thm:ca-schwarz-pick}
任意全纯映射 $f:\mathbb D\to\mathbb D$ 均满足
\[
 \left|\frac{f(z)-f(a)}{1-\overline{f(a)}f(z)}\right|
 \le\left|\frac{z-a}{1-\overline a z}\right|,
 \qquad
 \frac{|f'(a)|}{1-|f(a)|^2}\le\frac1{1-|a|^2}.
\]
若第一式在某对不同点处取等号，或第二式在某点取等号，则 $f$ 是圆盘自同构；反之，圆盘自同构处处保持等号。
\end{theorem}
\begin{proof}
令 $F=\varphi_{f(a)}\circ f\circ\varphi_a^{-1}$，则 $F:\mathbb D\to\mathbb D$ 全纯且 $F(0)=0$。施瓦茨引理的 $|F(\zeta)|\le|\zeta|$ 代入 $\zeta=\varphi_a(z)$ 就是第一式；其导数估计为
\[
 |F'(0)|=\frac{(1-|a|^2)|f'(a)|}{1-|f(a)|^2}\le1,
\]
给出第二式。上述任一取等号情形由施瓦茨引理使 $F$ 为旋转，解出 $f$ 为三个圆盘自同构的复合。反之，对自同构及其逆分别用不等式，即得等号。
\end{proof}

令 $\delta(z,a)=|(z-a)/(1-\overline a z)|$，称为圆盘的伪双曲距离（pseudohyperbolic distance）。上述结论说明 $\delta(f(z),f(a))\le\delta(z,a)$。对它作严格递增变换 $d(z,a)=2\operatorname{arctanh}\delta(z,a)$，也得到相同的压缩不等式；这给出庞加莱距离（Poincar\'e distance）的显式表达。这里使用的是已证的压缩关系，关于该距离的测地线等几何性质不参与后续证明。

\section{条带、扇形与裂平面的组合映射}

构造映射时，可把任务分成几个可验证步骤：先用平移、旋转确定位置，再用幂函数或指数改变开角和边界形状，最后用莫比乌斯变换送到圆盘或半平面。每一步都应注明定义域、分支、单射性与值域。

\begin{example}[水平条带映到单位圆盘]
设 $h>0$，求 $S=\{z:0<\operatorname{Im}z<h\}$ 到 $\mathbb D$ 的一个双全纯映射。
\end{example}
\begin{solve}
先取 $\zeta=e^{\pi z/h}$。其辐角在 $(0,\pi)$，模可取所有正数，所以像为上半平面。若两个条带内点有相同像，其差为 $2khi$；但虚部之差的绝对值小于 $h$，故 $k=0$，单射成立。反函数为 $(h/\pi)L(\zeta)$，其中选上半平面上辐角位于 $(0,\pi)$ 的对数。再接凯莱变换，得到
\[
 w=\frac{e^{\pi z/h}-i}{e^{\pi z/h}+i}.
\]
它把条带中点 $ih/2$ 映到原点。
\end{solve}

\begin{example}[扇形映到上半平面]
设 $0<\alpha<2\pi$，$S_\alpha=\{re^{i\theta}:r>0,\ 0<\theta<\alpha\}$。在这个扇形上选辐角 $\theta\in(0,\alpha)$，则
\[
 w=z^{\pi/\alpha}=\exp\left(\frac\pi\alpha(\ln r+i\theta)\right)
\]
把辐角区间变为 $(0,\pi)$，模 $r^{\pi/\alpha}$ 仍遍历所有正数，故双全纯映到上半平面。其逆变换使用上半平面上的辐角分支，再取 $\alpha/\pi$ 次幂。原点是扇形边界点，不在保角性要求的内部。
\end{example}

\begin{example}[裂平面映到圆盘]
对 $\Omega=\mathbb C\setminus(-\infty,0]$，先取主平方根 $s(z)=e^{\operatorname{Log}z/2}$。辐角范围从 $(-\pi,\pi)$ 变为 $(-\pi/2,\pi/2)$，故 $s$ 双全纯地映到右半平面，其逆为在右半平面上限制的平方函数。因此
\[
 w=\frac{\sqrt z-1}{\sqrt z+1}
\]
把 $\Omega$ 双全纯地映到单位圆盘，且把 $1$ 映到 $0$。这里 $\sqrt z$ 明确取主分支，反函数为 $z=((1+w)/(1-w))^2$。
\end{example}

\section{存在性与显式公式的边界}

\begin{example}[茹科夫斯基变换]
\emph{茹科夫斯基变换（Joukowski map）} $J(z)=\tfrac12(z+1/z)$ 将 $D=\{z:|z|>1\}$ 双全纯地映到 $\mathbb C\setminus[-1,1]$。
\end{example}
\begin{solve}
若 $J(z)=J(w)$，则 $(z-w)(1-1/(zw))=0$。在 $D$ 上 $|zw|>1$，故只能 $z=w$，得到单射。其导数 $J'(z)=(1-z^{-2})/2$ 在 $D$ 非零。

对任意 $\zeta\notin[-1,1]$，方程 $J(z)=\zeta$ 等价于 $z^2-2\zeta z+1=0$。两根之积为一；任何单位圆上的根 $e^{it}$ 都会使 $\zeta=\cos t\in[-1,1]$，因此此处两根都不在单位圆，恰有一个模大于一。反之，若 $\zeta\in[-1,1]$，可写成 $\cos t$，两根为 $e^{\pm it}$，均不在 $D$。这就确定了像集及满射性。单射全纯函数的局部逆可以拼成全局全纯逆，故结论成立。若把定义域扩大到整个去心平面，$J(z)=J(1/z)$ 通常产生两个不同原像，不能再声称单射。
\end{solve}

上面的例子依靠具体公式完成映射。对一般非空单连通真子域，是否总能映到圆盘，是更深的存在性问题。定理~\ref{thm:ca-riemann-mapping} 将在正规族工具建立后给出完整证明，并说明指定基点与正实导数如何保证唯一性。存在性不等于总能写出初等公式，也不自动保证任意粗糙边界上的连续一一延拓；这些差别将在该章分别讨论。

\section{习题}
\label{sec:ca-conformal-ex}
\begin{enumerate}
 \item 构造将 $0,1,\infty$ 依次映到 $1,\infty,0$ 的莫比乌斯变换。
 \item 构造上半平面到单位圆盘的双全纯映射，使 $2i$ 映到零。
 \item 构造第一象限到单位圆盘的双全纯映射，使 $1+i$ 映到零。
 \item 证明 $z^2$ 在右半平面单射，确定其像；解释它在整个去心平面为什么不单射。
 \item 若单位圆盘的双全纯自映射固定原点且导数为正实数，证明它是恒等映射。
\end{enumerate}

